% Livetime comparison presentation.
%
% Compile after making edits (run these commands from the repository root):
%   mkdir -p /tmp/livetime_slides_build
%   cd docs
%   pdflatex -interaction=nonstopmode -halt-on-error \
%     -output-directory=/tmp/livetime_slides_build livetime_analysis_update.tex
%   pdflatex -interaction=nonstopmode -halt-on-error \
%     -output-directory=/tmp/livetime_slides_build livetime_analysis_update.tex
%   cp /tmp/livetime_slides_build/livetime_analysis_update.pdf .

\documentclass[aspectratio=169,11pt]{beamer}

\usepackage[T1]{fontenc}
\usepackage{lmodern}
\usepackage{amsmath}
\usepackage{graphicx}
\usepackage{booktabs}

\definecolor{NPSBlue}{RGB}{18,62,98}
\definecolor{NPSTeal}{RGB}{0,139,139}
\definecolor{LightBlue}{RGB}{234,242,247}
\definecolor{DarkGray}{RGB}{50,55,60}

\setbeamercolor{normal text}{fg=DarkGray,bg=white}
\setbeamercolor{structure}{fg=NPSBlue}
\setbeamercolor{frametitle}{fg=white,bg=NPSBlue}
\setbeamercolor{block title}{fg=white,bg=NPSBlue}
\setbeamercolor{block body}{fg=DarkGray,bg=LightBlue}
\setbeamercolor{alerted text}{fg=NPSTeal}
\setbeamertemplate{navigation symbols}{}
\setbeamertemplate{blocks}[rounded][shadow=false]
\setbeamertemplate{itemize item}{\color{NPSTeal}$\blacktriangleright$}
\setbeamertemplate{itemize subitem}{\color{NPSTeal}$\bullet$}
\setbeamertemplate{footline}{%
  \hfill\raisebox{2mm}{\color{NPSBlue}\insertframenumber\hspace{3mm}}}
\setbeamersize{text margin left=7mm,text margin right=7mm}

\newcommand{\NEDTM}{N_{\mathrm{EDTM}}}
\newcommand{\Nhtrig}{N_{\mathrm{hTRIG4}}}
\newcommand{\beamOn}{f_{\mathrm{beam\ on}}}
\newcommand{\PlotPDF}{../output/efficiency_stuff/plots/compare_livetimes_multipanel_KinC_x60_4b_lh2.pdf}

\title{Livetime Definitions and EDTM Matching}
\subtitle{NPS livetime comparison}
\date{}

\begin{document}

% Slide 1
\begin{frame}[plain]
  \centering
  \vspace{1.5cm}
  {\color{NPSBlue}\LARGE\bfseries Livetime Definitions and EDTM Matching\par}
  \vspace{0.45cm}
  {\large NPS livetime comparison\par}
  \vspace{0.7cm}
  {\color{NPSTeal}\rule{0.45\textwidth}{1.5pt}}
\end{frame}

% Slide 2
\begin{frame}{Four definitions used in the NPS presentation}
\small
\begin{columns}[T,onlytextwidth]
  \begin{column}{0.49\textwidth}
    \begin{block}{1. Total EDTM livetime}
      \[
      LT_{\mathrm{EDTM}}=
      \frac{N_{\mathrm{EDTM}}^{\mathrm{TDC}}}
           {N_{\mathrm{EDTM}}^{\mathrm{scaler}}}
      \]
      \scriptsize Accepted: \(T.\mathrm{hEDTM}_{\rm raw}\neq0\). No current cut.
    \end{block}
    \begin{block}{3. Computer livetime all (TDC)}
      \[
      LT_{\mathrm{all,TDC}}=
      \frac{N_{\mathrm{hTRIG4}}^{\mathrm{TDC}}}
           {N_{\mathrm{hTRIG4}}^{\mathrm{scaler}}}\,\beamOn
      \]
      \scriptsize Accepted: \(T.\mathrm{hTRIG4}_{\rm raw}\neq0\).
    \end{block}
  \end{column}
  \begin{column}{0.49\textwidth}
    \begin{block}{2. Computer livetime all (TSH)}
      \[
      LT_{\mathrm{all,TSH}}=
      \frac{N_{\mathrm{hL1ACCP}}^{\mathrm{scaler}}}
           {N_{\mathrm{hTRIG4}}^{\mathrm{scaler}}}
      \]
      \scriptsize Accepted-trigger scaler: \(H.\mathrm{hL1ACCP}\). Current cut applied.
    \end{block}
    \begin{block}{4. Computer livetime physics (TDC)}
      \[
      LT_{\mathrm{phys,TDC}}=
      \frac{N_{\mathrm{physics}}^{\mathrm{TDC}}}
           {N_{\mathrm{hTRIG4}}^{\mathrm{scaler}}-N_{\mathrm{EDTM}}^{\mathrm{scaler}}}\,\beamOn
      \]
      \scriptsize Physics: hTRIG4 \(\neq0\) and EDTM \(=0\).
    \end{block}
  \end{column}
\end{columns}
\vspace{-1mm}
\centering\scriptsize
\(\displaystyle \beamOn=
N_{\mathrm{L1ACCP}}(\text{current cut})/N_{\mathrm{L1ACCP}}(\text{no cut})\).
\end{frame}

% Slide 3
\begin{frame}{Why preTRIG and hTRIG are different}
\centering
\vspace{2mm}
\begin{beamercolorbox}[rounded=true,sep=3mm,wd=0.94\textwidth]{block body}
  \centering\large
  Detector trigger hit
  \(\longrightarrow\)
  \alert{preTRIG\(_n\)}
  \(\longrightarrow\)
  Trigger Supervisor
  \(\longrightarrow\)
  \alert{hTRIG\(_n\)}
\end{beamercolorbox}

\vspace{5mm}
\begin{itemize}\setlength\itemsep{3mm}
  \item \textbf{preTRIG\(_n\)} fires for every hit delivered by trigger path \(n\).
  \item \textbf{hTRIG\(_n\)} is accumulated after further processing by the Trigger Supervisor.
  \item Trigger-Supervisor busy/acceptance therefore places an \textbf{implicit deadtime} in hTRIG\(_n\).
  \item hTRIG\(_n\) is consequently not the correct incident-trigger denominator. Use the matching pretrigger channel: \(\mathrm{PS}n\rightarrow\mathrm{preTRIG}n\), for \(n=1,\ldots,6\).
\end{itemize}
\end{frame}

% Slide 4
\begin{frame}{NewGen EDTM livetime}
\begin{columns}[T,onlytextwidth]
  \begin{column}{0.57\textwidth}
    \begin{block}{Exact common event selection}
    \small
    \begin{itemize}\setlength\itemsep{1.2mm}
      \item Run mean current, evaluated above \(2.5~\mu\mathrm{A}\), must be at least \(2.75~\mu\mathrm{A}\).
      \item Define \(I_0\) as the peak of the helicity-scaler current above \(2.5~\mu\mathrm{A}\); require \(0.85I_0\le I\le1.15I_0\).
      \item Require \(\texttt{actualHelicity}\neq0\) and retain complete helicity quartets.
      \item Over 30 scaler reads, require
      \(\left(Q_{\max}-Q_{\min}\right)/\overline Q\le0.08\).
      \item Require \(\texttt{fEvtHdr.fEvtType}=1\), \(g.\mathrm{evnum}\) inside the resulting accepted range, and event current inside the same current window.
    \end{itemize}
    \end{block}
  \end{column}
  \begin{column}{0.41\textwidth}
    \begin{block}{EDTM numerator}
    \small
    Require a finite raw EDTM value \(>1\), then count
    \[
      |t_{\rm raw}-t_{\rm raw}^{\rm peak}|\le500
      \quad\text{raw channels}. 
    \]
    The raw peak is the maximum of a 10-channel-binned histogram.
    \end{block}
    \begin{block}{Definition}
    \small
    \[
      LT_{\rm NewGen}=
      \mathrm{PS}\,
      \frac{N_{\rm raw\ peak}}
           {\Delta N_{\rm EDTM}^{\rm good\ current}}.
    \]
    \scriptsize The denominator includes every TSH interval whose current is inside the current window.
    \end{block}
  \end{column}
\end{columns}
\end{frame}

% Slide 5
\begin{frame}{Matched EDTM livetime}
\begin{columns}[T,onlytextwidth]
  \begin{column}{0.49\textwidth}
    \begin{block}{Event and interval selection}
    \footnotesize
    Start with the exact common selection from slide 4 and a finite raw EDTM value \(>1\).
    
    \vspace{2mm}
    Match the event by \(g.\mathrm{evnum}\) to a TSH interval
    \[
       [\,\mathrm{evNumber}_{i-1},\mathrm{evNumber}_{i}\,).
    \]
    Retain the interval only when it:
    \begin{itemize}\setlength\itemsep{1mm}
      \item lies wholly inside one accepted event-number range;
      \item has current in \([0.85I_0,1.15I_0]\);
      \item is fully covered by recorded events.
    \end{itemize}
    \end{block}
  \end{column}
  \begin{column}{0.49\textwidth}
    \begin{block}{Corrected-time match}
    \footnotesize
    \begin{enumerate}\setlength\itemsep{1mm}
      \item Use raw pulses within \(\pm500\) channels to form a corrected-time histogram with 0.1-ns bins.
      \item Find its rough maximum.
      \item Define the final peak as the median within \(\pm1\) ns of that maximum.
      \item Count an interval-matched pulse when
      \[
        |t_{\rm corr}-t_{\rm corr}^{\rm peak}|\le2~\mathrm{ns}.
      \]
    \end{enumerate}
    \end{block}
    \begin{block}{Definition}
    \small
    \[
      LT_{\rm matched}=
      \mathrm{PS}\,
      \frac{N_{\rm interval+time\ matched}}
           {\Delta N_{\rm EDTM}^{\rm selected\ intervals}}.
    \]
    \end{block}
  \end{column}
\end{columns}
\end{frame}

% Slide 6
\begin{frame}{Livetime comparison: KinC\_x60\_4b, LH2}
\vspace{-2mm}
\begin{columns}[c,onlytextwidth]
  \begin{column}{0.58\textwidth}
    \centering
    \includegraphics[page=1,width=\linewidth,height=0.62\textheight,keepaspectratio]{\PlotPDF}
  \end{column}
  \begin{column}{0.40\textwidth}
    \centering
    \includegraphics[page=2,width=\linewidth,height=0.28\textheight,keepaspectratio]{\PlotPDF}\\[1.5mm]
    \includegraphics[page=3,width=\linewidth,height=0.28\textheight,keepaspectratio]{\PlotPDF}
  \end{column}
\end{columns}
\end{frame}

% Slide 7
\begin{frame}{Summary}
\begin{itemize}\setlength\itemsep{4mm}
  \item The four original definitions have been implemented with the common good-event selection and the trigger channel selected from \(\mathrm{PS}1\)--\(\mathrm{PS}6\).
  \item hTRIG\(_n\) contains Trigger-Supervisor deadtime; trigger-based scaler denominators should use the corresponding preTRIG\(_n\).
  \item NewGen EDTM uses a broad raw-TDC peak and a good-current scaler denominator.
  \item Matched EDTM pairs events and scaler exposure in the same fully selected TSH intervals, then applies a \(\pm2\)-ns corrected-time window.
  \item The comparison output records every scaler quantity and plots all six definitions versus run number, beam current, and HMS S1X rate.
\end{itemize}
\end{frame}

\end{document}
